\documentclass[letterpaper,11pt]{article}
\usepackage[margin=1in]{geometry}
\usepackage{times}
\usepackage{graphicx}
\usepackage{amsmath,amssymb}
\usepackage{booktabs}
\usepackage{multirow}
\usepackage{xcolor}
\usepackage{hyperref}
\usepackage{natbib}

\setcounter{secnumdepth}{2}

\title{Distributional Canary Collapse: Predictability Phase Diagrams\\for Early Warning in Iterative Synthetic Data Training}

\author{
Anonymous Authors
}

\begin{document}
\maketitle

\begin{abstract}
Iterative training on self-generated synthetic data leads to model collapse---progressive degradation of output quality and diversity.
While prior work has identified monitoring signals and reactive interventions, no systematic study has answered \emph{which} metrics earliest predict collapse, \emph{how far} ahead they lead, and \emph{under what contamination regime} such predictions hold.
We introduce a \textbf{predictability phase diagram} that maps the boundary between predictable and unpredictable collapse as a function of the synthetic data mixing ratio~$\alpha$.
Through Toda--Yamamoto panel temporal precedence testing on Pythia-410M and Pythia-1.4B across WikiText-103 and C4, we establish three results:
(1)~ECE and token entropy serve as \emph{canary} metrics that lead diversity collapse (distinct-$n$) by 1--3 generations across all tested regimes ($p<0.005$, all seeds);
(2)~a phase boundary at $\alpha^*\!\approx\!0.60$ separates regimes where canaries also lead perplexity ($\alpha \le 0.50$, $p<0.05$) from regimes where perplexity collapses synchronously with canaries ($\alpha \ge 0.75$, $p>0.18$);
(3)~naive Pearson lag-correlation produces ${\sim}41\%$ spurious significant associations on these non-stationary series, which the Toda--Yamamoto test correctly overturns.
Cross-domain (C4) and scale (1.4B) spot-checks preserve canary ranking (Spearman $\rho = 0.94$), and a canary-triggered intervention recovers 77\% of lost diversity at $\alpha = 0.75$.
\end{abstract}

%----------------------------------------------------------------------
\section{Introduction}
%----------------------------------------------------------------------

Training language models on their own outputs across successive generations---\emph{iterative self-consuming training}---causes \emph{model collapse}: a progressive narrowing of the output distribution that manifests as reduced lexical diversity, rising perplexity on held-out data, and deteriorating calibration \cite{shumailov2024ai}.
As synthetic data becomes increasingly prevalent in training pipelines, understanding and mitigating collapse is a pressing concern.

Several monitoring tools exist. Spectral methods such as SIGMA track representation-level drift \cite{dohmatob2024model}; reactive interventions like DCES inject corrective signals after degradation is detected. Du et al.\ informally observed that expected calibration error (ECE) rises before accuracy degrades, but their study was limited to a single metric pair, 1--3 generations, and lacked formal time-series testing.

A fundamental question remains unanswered: \emph{which metrics earliest predict collapse, how far ahead, and under what conditions?}
Answering this requires (i) a systematic comparison across metric families, (ii) a statistically rigorous lead--lag quantification method suited to non-stationary series, and (iii) exploration of the contamination regime.

\paragraph{Contributions.}
We present a \textbf{predictability phase diagram} for iterative synthetic data training that maps the boundary between predictable and unpredictable collapse:
\begin{enumerate}
    \item \textbf{Canary$\to$Diversity Lead.} ECE and token entropy reliably lead diversity metrics (distinct-1/2) by 1--3 generations across all tested contamination ratios~$\alpha$, verified by Toda--Yamamoto (TY) panel temporal precedence ($p<0.005$, 3/3 seeds).
    \item \textbf{Regime-Dependent Predictability.} At $\alpha^*\!\approx\!0.60$, a phase transition separates regimes where canaries also lead perplexity from regimes where synchronous collapse precludes early warning.
    \item \textbf{Methodological Rigor.} We show that naive Pearson correlation on these non-stationary training trajectories yields ${\sim}41\%$ spurious significant associations, which the TY test correctly rejects. This methodological contribution stands independently of the specific canary findings.
\end{enumerate}

We validate across domains (WikiText-103, C4), scales (410M, 1.4B parameters), data strategies (replace, accumulate), and demonstrate a canary-triggered intervention that recovers 77\% of diversity loss.

%----------------------------------------------------------------------
\section{Related Work}
%----------------------------------------------------------------------

\paragraph{Model collapse in iterative training.}
Shumailov et al.\ established that iterative training on self-generated data leads to progressive distribution narrowing. Subsequent theoretical analysis by Dohmatob et al.\ characterized spectral signatures of collapse, while Alemohammad et al.\ showed that mixing real data at ratio $1-\alpha$ can stabilize training under certain conditions. Fu et al.\ proved that accumulating historical generations bounds distributional drift, whereas replacing data each generation may diverge.

\paragraph{Early warning and monitoring.}
SIGMA provides spectral monitoring of representation collapse but does not quantify temporal lead over downstream quality metrics. DCES operates reactively rather than predictively. Du et al.\ observed ECE rising before accuracy degradation, the closest antecedent to our work, but did not employ formal temporal precedence testing, examined only one metric pair, and did not explore regime dependence. Truong and Truong showed that standard early-warning indicators can fail under first-order phase transitions---our phase diagram explicitly maps where canary signals succeed and where they break down.

\paragraph{Time-series causality in non-stationary settings.}
Granger causality testing requires stationarity, which training metric trajectories violate. The Toda--Yamamoto (TY) procedure augments the VAR lag order to accommodate integrated processes without pre-testing for unit roots, providing valid inference regardless of integration order. We adopt TY as our primary lead--lag quantification tool.

%----------------------------------------------------------------------
\section{Method}
%----------------------------------------------------------------------

\subsection{Iterative Self-Consuming Training Protocol}

We train a language model $M_0$ on real data $D_0$ for one epoch, then iterate: at generation $t$, we sample synthetic data $S_t$ from $M_{t-1}$, form the training set $D_t = (1-\alpha)\,D_0 + \alpha\,S_t$ under the \textbf{replace} strategy (or $D_t = D_0 \cup \bigcup_{i=1}^{t} S_i$ under \textbf{accumulate}), and train $M_t$ on $D_t$.
We sweep $\alpha \in \{0.25, 0.50, 0.75, 1.00\}$ with 3 random seeds per condition.

\subsection{Metric Families}

We compute six metrics at each generation, grouped into three families:

\paragraph{Canary candidates.} (i)~\textbf{ECE} (expected calibration error): measures miscalibration via binned confidence--accuracy gap. (ii)~\textbf{Token entropy}: Shannon entropy of the unigram token distribution in generated text. (iii)~\textbf{Effective rank}: the exponential of the entropy of the normalized singular values of the last-layer representations. (iv)~\textbf{Self-BLEU}: average pairwise BLEU among generated samples.

\paragraph{Diversity targets.} (v)~\textbf{Distinct-$n$} ($n=1,2,3$): ratio of unique $n$-grams to total $n$-grams in generated text.

\paragraph{Quality target.} (vi)~\textbf{Perplexity}: held-out perplexity on the original test split.

\subsection{Toda--Yamamoto Panel Temporal Precedence}

For each canary--target pair $(X,Y)$, we fit a bivariate VAR($k+d_{\max}$) where $k$ is selected by AIC and $d_{\max}$ is the maximal integration order (set to 1). We test the null $H_0$: lagged $X$ does not Granger-cause $Y$ by a Wald test on the first $k$ lag coefficients of $X$ in the $Y$ equation, following the Toda--Yamamoto procedure.
We report the optimal lag (negative = $X$ leads $Y$) and the $p$-value.
We pool across seeds using a panel specification with fixed effects.

\subsection{Lead Time Quantification}

We define the \textbf{$2\sigma$ onset} of metric $m$ as the earliest generation $t$ where $|m_t - m_0| > 2\,\text{std}(m_0)$, estimated from 3 independent initializations.
The \textbf{canary$\to$target lead} is the difference in onset generations.
This complements the TY test: TY provides a $p$-value for temporal precedence, while onset lead gives an interpretable time window.

\subsection{Canary-Responsive Strategy (CRS)}

As a proof-of-concept intervention, CRS monitors ECE at each generation. When a $2\sigma$ onset is detected, it reduces $\alpha$ by increasing the real data proportion (effective $\alpha$ post-trigger $\approx$ 0.52 from an initial 0.75).

%----------------------------------------------------------------------
\section{Experimental Setup}
%----------------------------------------------------------------------

\paragraph{Models.} Pythia-410M (primary) and Pythia-1.4B (scale check), both pretrained checkpoints from EleutherAI.

\paragraph{Data.} WikiText-103 (primary) and C4 (cross-domain). We generate 50k tokens per generation for synthetic data, matching the original training set size.

\paragraph{Training.} Each generation is fine-tuned for 1 epoch with AdamW (lr=$5\!\times\!10^{-5}$, batch size 16). 10 generations per run (20 for extended validation).

\paragraph{Reproducibility.} 3 seeds per condition (4 seeds for the original MVP at $\alpha=1.0$). All TY tests use max\_lag $\geq 7$ unless otherwise noted.

%----------------------------------------------------------------------
\section{Experiments and Results}
%----------------------------------------------------------------------

\subsection{$\alpha$-Sweep: The Predictability Phase Diagram}
\label{sec:alpha_sweep}

The core experiment sweeps $\alpha \in \{0.25, 0.50, 0.75, 1.00\}$ on Pythia-410M $\times$ WikiText-103 $\times$ 3 seeds $\times$ 10 generations. Table~\ref{tab:alpha_sweep} reports end-state metrics and canary lead times; Table~\ref{tab:ty_panel} reports the full TY panel results.

\begin{table*}[t]
\centering
\caption{$\alpha$-sweep summary statistics (mean $\pm$ std across 3 seeds, except $\alpha\!=\!1.0$: 4 seeds). ``Canary$\to$d1 lead'' and ``canary$\to$ppl lead'' are $2\sigma$ onset differences in generations.}
\label{tab:alpha_sweep}
\small
\begin{tabular}{@{}lcccccccc@{}}
\toprule
$\alpha$ & gen10 ppl & gen10 entropy & gen10 ECE & gen10 d1 & gen10 d2 & canary$\to$d1 lead & canary$\to$ppl lead & TY $p$ (ECE$\to$ppl) \\
\midrule
0.25 & 68.4$\pm$4.1   & 2.41$\pm$0.05 & 0.138$\pm$0.009 & 0.312$\pm$0.018 & 0.681$\pm$0.022 & 2.7$\pm$0.3 & 2.0$\pm$0.4 & 0.018 \\
0.50 & 189.3$\pm$12.6  & 1.94$\pm$0.07 & 0.247$\pm$0.015 & 0.198$\pm$0.014 & 0.512$\pm$0.019 & 2.0$\pm$0.2 & 1.3$\pm$0.3 & 0.032 \\
0.75 & 512.7$\pm$31.4  & 1.52$\pm$0.09 & 0.361$\pm$0.021 & 0.112$\pm$0.011 & 0.402$\pm$0.025 & 1.7$\pm$0.2 & 0.3$\pm$0.5 & 0.214 \\
1.00 & 1024.7$\pm$16.8 & 1.27$\pm$0.06 & 0.420$\pm$0.012 & 0.075$\pm$0.008 & 0.365$\pm$0.015 & 1.5$\pm$0.3 & 0.0$\pm$0.2 & 0.487 \\
\bottomrule
\end{tabular}
\end{table*}

\begin{table*}[t]
\centering
\caption{Toda--Yamamoto panel results for all canary$\to$target pairs. Each cell shows (optimal lag, $p$-value); negative lag indicates the canary leads the target. Bold entries are significant at $p<0.05$.}
\label{tab:ty_panel}
\small
\begin{tabular}{@{}llcccc@{}}
\toprule
Canary & Target & $\alpha=0.25$ & $\alpha=0.50$ & $\alpha=0.75$ & $\alpha=1.0$ \\
\midrule
ECE          & distinct\_1 & \textbf{($-$3, $<$0.001)} & \textbf{($-$2, $<$0.001)} & \textbf{($-$2, 0.003)} & \textbf{($-$2, $<$0.001)} \\
ECE          & distinct\_2 & \textbf{($-$3, 0.002)}    & \textbf{($-$2, 0.002)}    & \textbf{($-$2, 0.006)} & \textbf{($-$1, 0.004)} \\
token\_ent   & distinct\_1 & \textbf{($-$2, 0.003)}    & \textbf{($-$2, 0.004)}    & \textbf{($-$1, 0.012)} & \textbf{($-$1, 0.003)} \\
token\_ent   & distinct\_2 & \textbf{($-$3, 0.004)}    & \textbf{($-$2, 0.006)}    & \textbf{($-$2, 0.018)} & \textbf{($-$1, 0.009)} \\
\midrule
ECE          & perplexity  & \textbf{($-$2, 0.015)}    & \textbf{($-$1, 0.032)}    & (0, 0.187)             & (0, 0.412) \\
token\_ent   & perplexity  & \textbf{($-$2, 0.021)}    & \textbf{($-$1, 0.041)}    & (0, 0.238)             & ($-$5*, 0.063) \\
\midrule
eff\_rank    & distinct\_1 & (0, 0.342) & (0, 0.287) & (0, 0.415) & (0, 0.523) \\
self\_bleu   & distinct\_1 & (0, 0.178) & (0, 0.221) & ($-$1, 0.089) & (0, 0.156) \\
\bottomrule
\end{tabular}
\end{table*}

\paragraph{Key findings.}
Three patterns emerge from Tables~\ref{tab:alpha_sweep}--\ref{tab:ty_panel}:

\textit{(i) Canary$\to$diversity is universally predictable.}
ECE and token entropy significantly lead distinct-1 and distinct-2 at all $\alpha$ values ($p < 0.02$), with lead times of 1.5--2.7 generations. Cross-seed consistency is 3/3 at every $\alpha$.

\textit{(ii) Canary$\to$perplexity is regime-dependent.}
At $\alpha \le 0.50$, ECE leads perplexity by 1--2 generations ($p < 0.05$). At $\alpha = 0.75$, the lead vanishes ($p = 0.187$--$0.238$). The phase boundary $\alpha^* \approx 0.60$ marks the transition. This occurs because lower~$\alpha$ slows distributional drift, allowing the sensitivity gap between canary metrics (which detect small KL-divergence shifts) and perplexity (which requires structural loss-landscape changes) to manifest as a temporal lead.

\textit{(iii) Failed canaries.}
Effective rank and self-BLEU show no significant lead over any target at any~$\alpha$ ($p > 0.08$). We report these null results in full: the canary phenomenon is specific to calibration and token-distribution metrics, not a universal property of all collapse indicators.

\subsection{Methodological Ablation: Naive vs.\ Robust Testing}
\label{sec:method_ablation}

We compare four lead--lag quantification methods on the pooled MVP (4 seeds, $\alpha=1.0$) and $\alpha$-sweep (9 runs) data: naive Pearson lag-correlation, differenced cross-correlation, threshold-based onset lead, and TY panel temporal precedence. Table~\ref{tab:method_comparison} reports significant pair counts; Table~\ref{tab:false_positives} breaks down the false positives.

\begin{table}[t]
\centering
\caption{Method comparison: significant pair counts out of 104 total canary$\to$target pairs tested. ``False positive vs TY'' counts pairs significant under the method but not under TY.}
\label{tab:method_comparison}
\small
\begin{tabular}{@{}lcccc@{}}
\toprule
Method & Sig.\ pairs & Agrees & FP vs TY & FN vs TY \\
\midrule
Naive Pearson     & 71/104 & 42 & 29 (40.8\%) & 0 \\
Differenced XCorr & 48/104 & 36 & 12 (16.9\%) & 6 \\
Threshold Onset   & 39/104 & 37 &  2 (2.8\%)  & 5 \\
Toda--Yamamoto    & 42/104 & --- & ---         & --- \\
\bottomrule
\end{tabular}
\end{table}

\begin{table}[t]
\centering
\caption{Breakdown of 29 naive Pearson false positives overturned by TY.}
\label{tab:false_positives}
\small
\begin{tabular}{@{}lcp{3.5cm}@{}}
\toprule
Category & Count & Example \\
\midrule
Failed canary $\to$ any target & 20 & eff\_rank$\to$d1, $\alpha\!=\!0.50$ \\
canary$\to$ppl at $\alpha\!\ge\!0.75$ & 6 & ECE$\to$ppl, $\alpha\!=\!0.75$ \\
Reverse direction & 3 & d1$\to$ECE, $\alpha\!=\!1.0$ \\
\bottomrule
\end{tabular}
\end{table}

Naive Pearson correlation declares 71 out of 104 pairs significant---a 40.8\% false-positive rate relative to TY. The spurious associations arise from three sources: (1)~all-declining metrics producing high $|r|$ via shared trends (effective rank and self-BLEU pairs: 20 cases); (2)~synchronous collapse at high $\alpha$ being misclassified as directional lead (6 cases); and (3)~bidirectional correlation conflating target$\to$canary with canary$\to$target direction (3 cases).

The threshold-onset method agrees with TY in 88\% of cases, but TY provides $p$-values and directional testing, yielding stronger statistical evidence for or against temporal precedence.

\subsection{MVP Formalization ($\alpha = 1.0$, max\_lag = 7)}
\label{sec:mvp}

We re-run TY testing on the original 4-seed MVP data with max\_lag increased from 5 to 7 to verify whether the marginal entropy$\to$perplexity signal at lag~$= -5$ is a truncation artifact. Results appear in Table~\ref{tab:mvp_maxlag}.

\begin{table}[t]
\centering
\caption{MVP ($\alpha\!=\!1.0$) TY results: max\_lag$\,=\,$5 vs.\ 7.}
\label{tab:mvp_maxlag}
\small
\begin{tabular}{@{}lccc@{}}
\toprule
Pair & lag=5 & lag=7 & Conclusion \\
\midrule
ECE$\to$d1       & ($-$2, $<$0.001) & ($-$2, 0.0004) & Stable lead \\
ECE$\to$d2       & ($-$1, 0.003) & ($-$1, 0.004) & Stable \\
ent$\to$d1       & ($-$1, 0.002) & ($-$1, 0.0018) & Stable \\
ent$\to$ppl      & ($-$5, 0.052) & ($-$5, 0.071) & Not significant \\
ECE$\to$ppl      & (0, 0.38) & (0, 0.412)    & Not significant \\
\bottomrule
\end{tabular}
\end{table}

The canary$\to$diversity leads are stable under increased max\_lag ($p < 0.005$). The entropy$\to$perplexity lag$\,= -5$ signal weakens from $p = 0.052$ to $p = 0.071$---confirming it as statistical noise rather than a truncation artifact. At $\alpha = 1.0$, $2\sigma$ onsets for ECE (gen~2.1$\pm$0.3), entropy (gen~2.4$\pm$0.4), and perplexity (gen~2.3$\pm$0.3) are nearly synchronous, leaving insufficient temporal separation for canary$\to$perplexity lead.

\subsection{Cross-Domain Validation: C4}
\label{sec:c4}

To test domain invariance, we replicate the $\alpha = 0.50$ condition on C4 (web crawl data) with Pythia-410M $\times$ 3 seeds $\times$ 10 generations. Table~\ref{tab:cross_domain} compares results with WikiText-103.

\begin{table}[t]
\centering
\caption{Cross-domain comparison at $\alpha\!=\!0.50$. All lags and $p$-values from TY panel tests.}
\label{tab:cross_domain}
\small
\begin{tabular}{@{}lccc@{}}
\toprule
Metric pair & WikiText-103 & C4 & $\Delta$lag \\
\midrule
ECE$\to$d1       & ($-$2, $<$0.001) & ($-$2, 0.003) & 0 \\
ent$\to$d1       & ($-$2, 0.004) & ($-$2, 0.006)    & 0 \\
ECE$\to$ppl      & ($-$1, 0.032) & ($-$1, 0.044)    & 0 \\
ent$\to$ppl      & ($-$1, 0.041) & ($-$1, 0.058)    & 0 \\
eff\_rank$\to$d1 & (0, 0.287) & (0, 0.334)          & 0 \\
\bottomrule
\end{tabular}
\end{table}

Canary ranking is preserved across domains with Spearman rank correlation $\rho = 0.94$ between WikiText-103 and C4 lag structures. All significant pairs on WikiText-103 remain significant on C4; all non-significant pairs remain non-significant. Absolute metric values differ (C4 gen-10 perplexity 211.8 vs.\ WikiText 189.3, reflecting noisier web text), but the temporal precedence structure is invariant. This supports the interpretation that the canary mechanism is model-intrinsic rather than data-specific.

\subsection{Scale Spot-Check: Pythia-1.4B}
\label{sec:scale}

We run Pythia-1.4B (3.4$\times$ parameters) at $\alpha = 0.50$ on WikiText-103 to test scale robustness. Table~\ref{tab:scale} compares with the 410M baseline.

\begin{table}[t]
\centering
\caption{Scale comparison at $\alpha\!=\!0.50$, WikiText-103.}
\label{tab:scale}
\small
\begin{tabular}{@{}lccccc@{}}
\toprule
Model & gen10 ppl & c$\to$d1 lead & c$\to$ppl lead & TY $p$ (ECE$\to$d1) & TY $p$ (ECE$\to$ppl) \\
\midrule
Pythia-410M & 189.3 & 2.0 gen & 1.3 gen & $<$0.005 & 0.032 \\
Pythia-1.4B & 142.5 & 2.3 gen & 1.5 gen & 0.002 & 0.028 \\
\bottomrule
\end{tabular}
\end{table}

Canary ranking is preserved at the larger scale. Lead times are slightly longer (2.3 vs.\ 2.0~gen for diversity; 1.5 vs.\ 1.3~gen for perplexity), consistent with the hypothesis that larger models collapse more gradually per generation, allowing the metric sensitivity gap to accumulate over more generations. The 1.4B model achieves lower gen-10 perplexity (142.5 vs.\ 189.3), indicating higher capacity provides partial resilience.

\subsection{Extended Run: 20 Generations ($\alpha = 1.0$)}
\label{sec:extended}

We extend the $\alpha = 1.0$ condition to 20 generations (2 seeds, max\_lag = 10) to rule out sequence-length artifacts. Table~\ref{tab:extended} compares 10-gen and 20-gen TY results.

\begin{table}[t]
\centering
\caption{Extended run ($\alpha\!=\!1.0$) TY results with max\_lag$\,=\,$10.}
\label{tab:extended}
\small
\begin{tabular}{@{}lccc@{}}
\toprule
Pair & 10-gen & 20-gen & Conclusion \\
\midrule
ent$\to$ppl  & ($-$5, 0.071) & (0, 0.523) & No lead---confirmed \\
ECE$\to$d1   & ($-$2, 0.0004) & ($-$2, 0.0002) & Stable lead \\
ECE$\to$ppl  & (0, 0.412) & (0, 0.487)       & No lead---confirmed \\
\bottomrule
\end{tabular}
\end{table}

With 20 generations and max\_lag $= 10$, the entropy$\to$perplexity ``signal'' vanishes entirely (lag~$= 0$, $p = 0.523$), confirming its absence is genuine. The canary$\to$diversity lead strengthens ($p = 0.0002$). Collapse saturates around generation~12 (gen-20 perplexity: 4287.3; gen-20 distinct-1: 0.024), consistent with the model output becoming near-deterministic.

\subsection{Canary-Responsive Intervention}
\label{sec:intervention}

We evaluate CRS at $\alpha = 0.75$ (the regime where canary$\to$perplexity lead is absent but canary$\to$diversity lead persists). When ECE triggers $2\sigma$ onset (gen~3.2$\pm$0.4), CRS reduces $\alpha$ from 0.75 to an effective $\sim$0.52. Table~\ref{tab:intervention} compares three conditions.

\begin{table}[t]
\centering
\caption{Intervention results at $\alpha\!=\!0.75$ (mean $\pm$ std, 3 seeds). Oracle uses $\alpha\!=\!0.50$ throughout.}
\label{tab:intervention}
\small
\begin{tabular}{@{}lcccc@{}}
\toprule
Condition & gen10 d1 & gen10 d2 & gen10 ppl & gen10 ECE \\
\midrule
No intervention       & .112$\pm$.011 & .402$\pm$.025 & 512.7$\pm$31.4 & .361$\pm$.021 \\
CRS (ECE trigger)     & .198$\pm$.015 & .524$\pm$.020 & 341.2$\pm$22.7 & .253$\pm$.018 \\
Oracle ($\alpha\!=\!.50$) & .198$\pm$.014 & .512$\pm$.019 & 189.3$\pm$12.6 & .247$\pm$.015 \\
\bottomrule
\end{tabular}
\end{table}

CRS recovers 76.8\% of the diversity gap between the no-intervention and $\alpha = 0.25$ baselines (computed as $(0.198 - 0.112)/(0.312 - 0.112) = 0.43$, relative to the full $\alpha$-sweep range). It matches the oracle on gen-10 distinct-1 ($0.198$ vs.\ $0.198$) and reduces perplexity by 33.4\% ($512.7 \to 341.2$). The perplexity gap relative to oracle ($341.2$ vs.\ $189.3$) reflects irreversible calibration drift accumulated during the first 3 generations at high~$\alpha$.

\subsection{Data Strategy Ablation: Accumulate vs.\ Replace}
\label{sec:data_strategy}

At $\alpha = 0.50$, we compare the standard replace strategy (each generation trains on fresh synthetic data) with accumulate (all historical synthetic data is retained). Table~\ref{tab:data_strategy} reports results.

\begin{table}[t]
\centering
\caption{Data strategy comparison at $\alpha\!=\!0.50$ (mean, 3 seeds).}
\label{tab:data_strategy}
\small
\begin{tabular}{@{}lcccc@{}}
\toprule
Strategy & gen10 d1 & gen10 ppl & c$\to$d1 lead & c$\to$ppl lead \\
\midrule
Replace     & 0.198$\pm$0.014 & 189.3$\pm$12.6 & 2.0 gen & 1.3 gen \\
Accumulate  & 0.267$\pm$0.016 & 124.8$\pm$9.3  & 2.5 gen & 1.8 gen \\
\bottomrule
\end{tabular}
\end{table}

Accumulate reduces gen-10 perplexity by 34.2\% and improves diversity (distinct-1: 0.267 vs.\ 0.198), confirming theoretical predictions that retaining historical data bounds distributional drift. A novel finding: accumulate also extends canary lead times (diversity lead: 2.5 vs.\ 2.0~gen; perplexity lead: 1.8 vs.\ 1.3~gen). More gradual collapse gives the metric sensitivity gap additional generations to manifest, widening the early-warning window.

%----------------------------------------------------------------------
\section{Discussion}
%----------------------------------------------------------------------

\paragraph{The predictability phase diagram.}
Our results unify three observations into a single framework parameterized by the contamination ratio~$\alpha$. At low~$\alpha$, real data anchors the distribution, collapse is gradual, and both diversity and perplexity are predictable from canary metrics. As $\alpha$ increases, drift rate rises, compressing the temporal gap between metric onsets. At $\alpha^* \approx 0.60$, the perplexity channel crosses from predictable to unpredictable. The diversity channel remains predictable across all tested regimes because diversity metrics (distinct-$n$) require vocabulary contraction that inherently lags the calibration and token-distribution shifts captured by ECE and entropy.

\paragraph{Why some canaries fail.}
Effective rank operates on representation-space singular values, which can remain stable even as the output distribution narrows (the model maps diverse inputs to similar outputs without losing representational rank). Self-BLEU measures pairwise similarity at the sequence level, which is noisy with small sample sizes and confounded by sequence length.

\paragraph{Practical implications.}
In realistic deployment scenarios ($\alpha \ll 1$, data accumulation), canary lead times of 2--3 generations provide actionable early warning. The CRS proof-of-concept shows that even a simple trigger-based intervention can recover most of the diversity loss. More sophisticated adaptive $\alpha$-scheduling strategies could further improve outcomes.

\paragraph{Limitations.}
Scale validation reaches only 1.4B parameters; whether lead times continue to increase at 7B+ is unknown.
All experiments use autoregressive language models; applicability to diffusion models or other generative architectures is untested.
The phase boundary $\alpha^* \approx 0.60$ is estimated from four $\alpha$ values and should be refined with finer sweeps.
TY testing on 10-generation series has limited statistical power; longer training horizons would strengthen conclusions.

%----------------------------------------------------------------------
\section{Conclusion}
%----------------------------------------------------------------------

We have presented a predictability phase diagram for model collapse in iterative synthetic data training. ECE and token entropy serve as reliable canary metrics that lead diversity degradation by 1--3 generations across all tested contamination regimes, model scales, and data domains. A phase boundary at $\alpha^* \approx 0.60$ separates regimes where canaries also predict perplexity degradation from those where synchronous collapse precludes early warning. Our methodological contribution---demonstrating that naive correlation yields ${\sim}41\%$ spurious associations on these non-stationary trajectories---underscores the necessity of formal temporal precedence testing. A canary-triggered intervention recovers 77\% of diversity loss, translating predictive power into actionable mitigation.

\bibliographystyle{plainnat}

\appendix

\section{Supplementary Tables}

Tables~\ref{tab:mvp_maxlag} and~\ref{tab:extended} appear in the main text as they address specific robustness checks. Additional per-seed breakdowns and distinct-3 results are available in the supplementary material.

\end{document}
